%%%%%%%%%%%%%%%%%%%%%%%%%%%%%%%%%%%%%%%%%%%%%%%%%%%%%%%%%%%%%%%%%%%%%%
%% Copyright (c) 2016 Carsten Wulff Software, Norway
%% %%%%%%%%%%%%%%%%%%%%%%%%%%%%%%%%%%%%%%%%%%%%%%%%%%%%%%%%%%%%%%%%%%%
%% Created       : wulff at 2016-11-5
%% %%%%%%%%%%%%%%%%%%%%%%%%%%%%%%%%%%%%%%%%%%%%%%%%%%%%%%%%%%%%%%%%%%%
%% This program is free software: you can redistribute it and/or modify
%% it under the terms of the GNU General Public License as published by
%% the Free Software Foundation, either version 3 of the License, or
%% (at your option) any later version.
%%
%% This program is distributed in the hope that it will be useful,
%% but WITHOUT ANY WARRANTY; without even the implied warranty of
%% MERCHANTABILITY or FITNESS FOR A PARTICULAR PURPOSE.  See the
%% GNU General Public License for more details.
%%
%% You should have received a copy of the GNU General Public License
%% along with this program.  If not, see <http://www.gnu.org/licenses/>.
%%%%%%%%%%%%%%%%%%%%%%%%%%%%%%%%%%%%%%%%%%%%%%%%%%%%%%%%%%%%%%%%%%%%%%
%\documentclass[final,11pt]{IEEEtran}
\documentclass[crop,border=0pt]{standalone}
\usepackage{times}


%\usepackage{placeins}
\usepackage{upgreek}
\usepackage{amssymb,amsmath}
\usepackage{cite}
\usepackage{verbatim}
\usepackage{listings}
\usepackage{xcolor}
\usepackage{flafter}
\usepackage{booktabs}
\usepackage{url}
\usepackage{ textcomp }

\hyphenation{op-tical net-works semi-conduc-tor}


% -JSON styel

%\definecolor{darkGreen}{rgb}{0,0.2097,0}
%\newcommand{\myfigwidth}{\linewidth}
%\newcommand{\myfigwidtha}{\linewidth}
%\newcommand{\myfigwidthl}{\textwidth}
%\newcommand{\myfigwidthc}{0.6\linewidth}


\newcommand{\myfigwidtha}{3.3in}
\newcommand{\myfigwidthl}{7in}
\newcommand{\myfigwidthc}{2in}
%\newcommand{\myfigwidthe}{3.2in}
%\newcommand{\myfigwidthb}{5in}

\newcommand{\myfigname}{fig_}
\newcommand{\reg}[2]{{\figurename} \ref{#1}(#2)}
\newcommand{\req}[1]{{\figurename} \ref{#1}}
%\newcommand{\myeqname}{eq:}
%\newcommand{\req}[1]{(\ref{\myeqname#1})}

\definecolor{armygreen}{rgb}{0, 0.5, 0}
\newcommand{\missing}[1]{{ #1}}
\newcommand{\edit}[1]{{ #1}}
\newcommand{\secedit}[1]{{ #1}}
\newcommand{\deleted}[1]{{ #1}}
\newcommand{\moved}[1]{{ #1}}

%\newcommand{\missing}[1]{{\color{red} #1}}
%\newcommand{\edit}[1]{{\color{red} #1}}
%\newcommand{\secedit}[1]{{\color{armygreen} #1}}
%\newcommand{\deleted}[1]{{\color{violet} #1}}
%\newcommand{\moved}[1]{{\color{blue} #1}}

\newcommand{\SD}{Sigma-Delta }
\newcommand{\eqn}[1]{
  \begin{math}
    #1
  \end{math}}

\newcommand\JSONnumbervaluestyle{\color{blue}}
\newcommand\JSONstringvaluestyle{\color{red}}



\newif\ifcolonfoundonthisline

\makeatletter

\lstdefinestyle{json}
{
  showstringspaces    = false,
  keywords            = {false,true},
  alsoletter          = 0123456789.,
  morestring          = [s]{"}{"},
  stringstyle         = \ifcolonfoundonthisline\JSONstringvaluestyle\fi,
  MoreSelectCharTable =%
  \lst@DefSaveDef{`:}\colon@json{\processColon@json},
  basicstyle          = \fontsize{7}{7}\ttfamily,
  keywordstyle        = \fontsize{7}{7}\ttfamily\bfseries,
}

% flip the switch if a colon is found in Pmode
\newcommand\processColon@json{%
  \colon@json%
  \ifnum\lst@mode=\lst@Pmode%
  \global\colonfoundonthislinetrue%
  \fi
}

\lst@AddToHook{Output}{%
  \ifcolonfoundonthisline%
  \ifnum\lst@mode=\lst@Pmode%
  \def\lst@thestyle{\JSONnumbervaluestyle}%
  \fi
  \fi
  % override by keyword style if a keyword is detected!
  \lsthk@DetectKeywords%
}

% reset the switch at the end of line
\lst@AddToHook{EOL}%
{\global\colonfoundonthislinefalse}

\makeatother

% \lstset{  basicstyle=\tiny}

% \lstset{
% basicstyle=\fontsize{11}{13}\selectfont\ttfamily
% }


\lstset{language=Java}




%\usepackage[paperwidth=\maxdimen,paperheight=\maxdimen]{geometry}


\usepackage{tikz,pgfplots}
\usepackage{ifthen}
\usepackage{calc}
\usetikzlibrary{patterns}
%\usepackage[tightpage,active]{preview}
\usetikzlibrary{circuits.logic.US}
\usetikzlibrary{circuits.ee.IEC}
\usepgflibrary{shapes.arrows}
\usepgflibrary{arrows}
\usepackage{circuitikz}


% circuitikz draws bipole outlines (resistors, diodes, capacitors) at
% bipoles/thickness times the ambient line width, and the default is 2.
% That makes a circuitikz resistor visibly heavier than the wire it sits
% on, and heavier than the hand rolled zig-zag in ckt_lib.tex. One is the
% house width: everything is drawn at the environment's own thickness.
% Arrow tips: one filled tip everywhere, set through the '>' knob so a
% figure only ever writes ->, <- or <->. Latex (from arrows.meta)
% scales with the line width, so it stays in proportion if a drawing
% is ever set thicker or thinner than the house 'thick'.
\usetikzlibrary{arrows.meta}
\tikzset{>={Latex}}

\ctikzset{bipoles/thickness=1}
\ctikzset{tripoles/mos style/arrows}
\ctikzset{tripoles/pmos style/emptycircle}
%\PreviewEnvironment{circuitikz}
\pagestyle{empty}


\tikzstyle{cicfill}=[minimum size=1.4em]
\tikzstyle{stuff_fill}=[rectangle,fill=white,minimum size=1.4em]

\newcommand{\grayOne}{0.3 }
\newcommand{\grayTwo}{0.15 }
\newcommand{\grayThree}{0.26 }
\newcommand{\grayFour}{0.41 }
\newcommand{\grayFive}{0.61 }
\newcommand{\graySix}{0.85 }

%- Gray tone
%\definecolor{active}{rgb}{\graySix,\graySix,\graySix}
%\definecolor{poly}{rgb}{\grayFive,\grayFive,\grayFive}
%\definecolor{contact}{rgb}{\graySix,\graySix,\graySix}
%\definecolor{mOne}{rgb}{\grayFour,\grayFour,\grayFour}
%\definecolor{mTwo}{rgb}{\graySix,\graySix,\graySix}
%\definecolor{mThree}{rgb}{\grayThree,\grayThree,\grayThree}
%\definecolor{mFour}{rgb}{\grayTwo,\grayTwo,\grayTwo}

\definecolor{grey}{rgb}{1,1,1}
\definecolor{lightgrey}{rgb}{\grayFour,\grayFour,\grayFour}
\definecolor{active}{rgb}{0.4,1,0.6}
\definecolor{poly}{rgb}{1.0,0.5,0.5}
\definecolor{cut}{rgb}{0.91,0.91,0}
\definecolor{mOne}{rgb}{0.36,0.36,1}
\definecolor{mTwo}{rgb}{0.655,0.655,0}
\definecolor{mThree}{rgb}{0.3,1,1}
\definecolor{mFour}{rgb}{0,0.2097,0}
\begin{document}


% In-band SNDR and SFDR of the 1-bit modulator against input level, OSR = 1024.
%
% A longer record than the spectra above, 262144 samples with the input
% at bin 37, so the band holds enough bins to measure in. Dither
% power is relative to the quantization noise power, step^2/12. The
% dashed line is the first order SQNR for the actual step. Without
% dither the loop tracks it near full scale, where a large input dithers the loop by
% itself, and falls away from it as the input shrinks, until below about
% -68 dBFS the output is a limit cycle that does not see the input at
% all. With equal-power dither the SNDR follows the estimate about 4 dB
% below it, and below -40 dBFS the SFDR is 18 to 35 dB
% better. Dither is what makes the white-noise model true; near
% full scale it costs SNDR.
%
% Generated by py/tikzplot.py. Edit the script in ex/ that
% produced it and re-run that, not this file.

\begin{tikzpicture}[thick]

\begin{groupplot}[
  group style={group size=2 by 1, horizontal sep=1.7cm, vertical sep=1.5cm},
  scale only axis,
  grid=both,
  grid style={gray!22, very thin},
  tick align=outside,
  tick label style={font=\small},
  label style={font=\small},
  title style={font=\small, yshift=-1mm},
  legend cell align=left,
  legend style={font=\small, draw=gray!50, fill=white, fill opacity=0.85, text opacity=1}
]
\nextgroupplot[
  width=6.0cm,
  height=5.5cm,
  xlabel={Input amplitude [dBFS]},
  ylabel={[dB]},
  title={In-band SNDR},
  ymin=-40,
  ymax=100,
  legend pos=south east
]
\addplot[black, very thick, mark=none] coordinates {(-80,-13.376) (-76,-13.376) (-72,-13.376) (-68,-13.376) (-64,1.1255) (-60,6.6481) (-56,11.734) (-52,16.992) (-48,19.261) (-44,30.035) (-40,35.652) (-36,40.193) (-32,47.238) (-28,51.942) (-24,60.039) (-20,65.757) (-16,71.856) (-12,76.885) (-8,80.107) (-4,85.394) (0,78.467)};
\addlegendentry{no dither}
\addplot[blue, very thick, mark=none] coordinates {(-80,9.0702) (-76,4.1998) (-72,-0.74215) (-68,-1.6892) (-64,10.783) (-60,13.008) (-56,23.58) (-52,26.996) (-48,29.865) (-44,38.677) (-40,45.236) (-36,47.629) (-32,54.919) (-28,57.851) (-24,64.014) (-20,68.867) (-16,73.197) (-12,77.587) (-8,81.335) (-4,85.445) (0,72.916)};
\addlegendentry{$-12$ dB dither}
\addplot[red, very thick, mark=none] coordinates {(-80,3.3537) (-76,7.276) (-72,11.239) (-68,15.238) (-64,19.387) (-60,23.367) (-56,27.114) (-52,30.944) (-48,35.472) (-44,39.235) (-40,43.31) (-36,46.95) (-32,50.864) (-28,55.067) (-24,59.181) (-20,63.089) (-16,67.11) (-12,70.795) (-8,74.889) (-4,76.404) (0,61.473)};
\addlegendentry{0 dB dither}
\addplot[black, thick, dashed, mark=none] coordinates {(-80,6.8999) (-76,10.9) (-72,14.9) (-68,18.9) (-64,22.9) (-60,26.9) (-56,30.9) (-52,34.9) (-48,38.9) (-44,42.9) (-40,46.9) (-36,50.9) (-32,54.9) (-28,58.9) (-24,62.9) (-20,66.9) (-16,70.9) (-12,74.9) (-8,78.9) (-4,82.9) (0,86.9)};
\addlegendentry{SQNR estimate}
\nextgroupplot[
  width=6.0cm,
  height=5.5cm,
  xlabel={Input amplitude [dBFS]},
  title={In-band SFDR},
  ymin=-40,
  ymax=100
]
\addplot[black, very thick, mark=none] coordinates {(-80,-9.5424) (-76,-9.5424) (-72,-9.5424) (-68,-9.5424) (-64,1.8322) (-60,6.6742) (-56,14.107) (-52,18.268) (-48,19.265) (-44,31.089) (-40,38.656) (-36,40.922) (-32,48.979) (-28,52.419) (-24,61.206) (-20,68.729) (-16,73.716) (-12,79.168) (-8,80.875) (-4,86.251) (0,81.07)};
\addplot[blue, very thick, mark=none] coordinates {(-80,11.148) (-76,4.5463) (-72,-0.034089) (-68,-0.78452) (-64,12.087) (-60,15.423) (-56,26.597) (-52,27.701) (-48,30.184) (-44,39.673) (-40,47.663) (-36,48.69) (-32,58.977) (-28,59.803) (-24,68.343) (-20,75.963) (-16,80.554) (-12,86.602) (-8,88.925) (-4,93.522) (0,74.373)};
\addplot[red, very thick, mark=none] coordinates {(-80,16.87) (-76,20.707) (-72,25.144) (-68,28.498) (-64,32.784) (-60,36.3) (-56,40.967) (-52,43.857) (-48,49.368) (-44,52.091) (-40,56.842) (-36,60.717) (-32,64.435) (-28,68.823) (-24,72.419) (-20,76.363) (-16,81.139) (-12,84.262) (-8,87.453) (-4,81.223) (0,61.757)};
\end{groupplot}

\end{tikzpicture}
\end{document}
